    
    

    
    \hypertarget{a-minimal-alf-run}{%
\section{A minimal ALF run}\label{a-minimal-alf-run}}

    In this bare-bones example we use the
\href{https://git.physik.uni-wuerzburg.de/ALF/pyALF/-/tree/master/}{pyALF}
interface to run the canonical Hubbard model on a default configuration:
a \(6\times6\) square grid, with interaction strength \(U=4\) and
inverse temperature \(\beta = 5\).

Bellow we go through the steps for performing the simulation and
outputting observables.

\begin{center}\rule{0.5\linewidth}{0.5pt}\end{center}

\textbf{1.} Import \texttt{ALF\_source} and \texttt{Simulation} classes
from the \texttt{py\_alf} python module, which provide the interface
with ALF:

    \begin{tcolorbox}[breakable, size=fbox, boxrule=1pt, pad at break*=1mm,colback=cellbackground, colframe=cellborder]
\prompt{In}{incolor}{1}{\boxspacing}
\begin{Verbatim}[commandchars=\\\{\}]
\PY{k+kn}{from} \PY{n+nn}{py\PYZus{}alf} \PY{k+kn}{import} \PY{n}{ALF\PYZus{}source}\PY{p}{,} \PY{n}{Simulation}  \PY{c+c1}{\PYZsh{} Interface with ALF}
\end{Verbatim}
\end{tcolorbox}

    \textbf{2.} Create an instance of \texttt{ALF\_source}, downloading the
ALF source code from the
\href{https://git.physik.uni-wuerzburg.de/ALF/ALF/-/tree/master/}{ALF
repository}, if \texttt{alf\_dir} does not exist. Gets alf\_dir from
environment variable \texttt{\$ALF\_DIR}, or defaults to ``./ALF'', if
not present:

    \begin{tcolorbox}[breakable, size=fbox, boxrule=1pt, pad at break*=1mm,colback=cellbackground, colframe=cellborder]
\prompt{In}{incolor}{2}{\boxspacing}
\begin{Verbatim}[commandchars=\\\{\}]
\PY{n}{alf\PYZus{}src} \PY{o}{=} \PY{n}{ALF\PYZus{}source}\PY{p}{(}\PY{p}{)}
\end{Verbatim}
\end{tcolorbox}

    \textbf{3.} Create an instance of \texttt{Simulation}, overwriting
default parameters as desired:

    \begin{tcolorbox}[breakable, size=fbox, boxrule=1pt, pad at break*=1mm,colback=cellbackground, colframe=cellborder]
\prompt{In}{incolor}{3}{\boxspacing}
\begin{Verbatim}[commandchars=\\\{\}]
\PY{n}{sim} \PY{o}{=} \PY{n}{Simulation}\PY{p}{(}
    \PY{n}{alf\PYZus{}src}\PY{p}{,}
    \PY{l+s+s2}{\PYZdq{}}\PY{l+s+s2}{Hubbard}\PY{l+s+s2}{\PYZdq{}}\PY{p}{,}                    \PY{c+c1}{\PYZsh{} Name of Hamiltonian}
    \PY{p}{\PYZob{}}                             \PY{c+c1}{\PYZsh{} Dictionary overwriting default parameters}
        \PY{l+s+s2}{\PYZdq{}}\PY{l+s+s2}{Lattice\PYZus{}type}\PY{l+s+s2}{\PYZdq{}}\PY{p}{:} \PY{l+s+s2}{\PYZdq{}}\PY{l+s+s2}{Square}\PY{l+s+s2}{\PYZdq{}}
    \PY{p}{\PYZcb{}}\PY{p}{,}
    \PY{n}{machine}\PY{o}{=}\PY{l+s+s1}{\PYZsq{}}\PY{l+s+s1}{GNU}\PY{l+s+s1}{\PYZsq{}}  \PY{c+c1}{\PYZsh{} Change to \PYZdq{}intel\PYZdq{}, or \PYZdq{}PGI\PYZdq{} if gfortran is not installed}
\PY{p}{)}
\end{Verbatim}
\end{tcolorbox}

    \textbf{4.} Compile ALF. The first time it will also download and
compile HDF5, which could take \(\sim15\) minutes.

    \begin{tcolorbox}[breakable, size=fbox, boxrule=1pt, pad at break*=1mm,colback=cellbackground, colframe=cellborder]
\prompt{In}{incolor}{4}{\boxspacing}
\begin{Verbatim}[commandchars=\\\{\}]
\PY{n}{sim}\PY{o}{.}\PY{n}{compile}\PY{p}{(}\PY{p}{)}
\end{Verbatim}
\end{tcolorbox}

    \begin{Verbatim}[commandchars=\\\{\}]
Compiling ALF{\ldots}
Done.
    \end{Verbatim}

    \textbf{5.} Perform the simulation as specified in \texttt{sim}:

    \begin{tcolorbox}[breakable, size=fbox, boxrule=1pt, pad at break*=1mm,colback=cellbackground, colframe=cellborder]
\prompt{In}{incolor}{5}{\boxspacing}
\begin{Verbatim}[commandchars=\\\{\}]
\PY{n}{sim}\PY{o}{.}\PY{n}{run}\PY{p}{(}\PY{p}{)}
\end{Verbatim}
\end{tcolorbox}

    \begin{Verbatim}[commandchars=\\\{\}]
Prepare directory "/home/stafusa/ALF/pyALF/Notebooks/ALF\_data/Hubbard\_Square"
for Monte Carlo run.
Create new directory.
Run /home/stafusa/ALF/ALF/Prog/ALF.out
    \end{Verbatim}

    \textbf{6.} Perform some simple analysis:

    \begin{tcolorbox}[breakable, size=fbox, boxrule=1pt, pad at break*=1mm,colback=cellbackground, colframe=cellborder]
\prompt{In}{incolor}{6}{\boxspacing}
\begin{Verbatim}[commandchars=\\\{\}]
\PY{n}{sim}\PY{o}{.}\PY{n}{analysis}\PY{p}{(}\PY{p}{)}
\end{Verbatim}
\end{tcolorbox}

    \begin{Verbatim}[commandchars=\\\{\}]
\#\#\# Analyzing /home/stafusa/ALF/pyALF/Notebooks/ALF\_data/Hubbard\_Square \#\#\#
/home/stafusa/ALF/pyALF/Notebooks
Custom observables:
Scalar observables:
Ener\_scal
Kin\_scal
Part\_scal
Pot\_scal
Histogram observables:
Equal time observables:
Den\_eq
Green\_eq
SpinT\_eq
SpinXY\_eq
SpinZ\_eq
Time displaced observables:
Den\_tau
Green\_tau
SpinT\_tau
SpinXY\_tau
SpinZ\_tau
    \end{Verbatim}

    \textbf{7.} Read analysis results into a Pandas Dataframe with one row
per simulation, containing parameters and observables:

    \begin{tcolorbox}[breakable, size=fbox, boxrule=1pt, pad at break*=1mm,colback=cellbackground, colframe=cellborder]
\prompt{In}{incolor}{7}{\boxspacing}
\begin{Verbatim}[commandchars=\\\{\}]
\PY{n}{obs} \PY{o}{=} \PY{n}{sim}\PY{o}{.}\PY{n}{get\PYZus{}obs}\PY{p}{(}\PY{p}{)}
\end{Verbatim}
\end{tcolorbox}

    \begin{Verbatim}[commandchars=\\\{\}]
/home/stafusa/ALF/pyALF/Notebooks/ALF\_data/Hubbard\_Square
    \end{Verbatim}

    \begin{tcolorbox}[breakable, size=fbox, boxrule=1pt, pad at break*=1mm,colback=cellbackground, colframe=cellborder]
\prompt{In}{incolor}{8}{\boxspacing}
\begin{Verbatim}[commandchars=\\\{\}]
\PY{n}{obs}
\end{Verbatim}
\end{tcolorbox}

            \begin{tcolorbox}[breakable, size=fbox, boxrule=.5pt, pad at break*=1mm, opacityfill=0]
\prompt{Out}{outcolor}{8}{\boxspacing}
\begin{Verbatim}[commandchars=\\\{\}]
                                                    ham\_t  ham\_chem  ham\_u  \textbackslash{}
/home/stafusa/ALF/pyALF/Notebooks/ALF\_data/Hubb{\ldots}    1.0       0.0    0.1

                                                    ham\_t2  ham\_u2  ham\_tperp  \textbackslash{}
/home/stafusa/ALF/pyALF/Notebooks/ALF\_data/Hubb{\ldots}     1.0     0.1        1.0

                                                    mz  continuous  \textbackslash{}
/home/stafusa/ALF/pyALF/Notebooks/ALF\_data/Hubb{\ldots}   1           0

                                                         model lattice\_type  \textbackslash{}
/home/stafusa/ALF/pyALF/Notebooks/ALF\_data/Hubb{\ldots}  b'Hubbard'    b'Square'

                                                    {\ldots}  \textbackslash{}
/home/stafusa/ALF/pyALF/Notebooks/ALF\_data/Hubb{\ldots}  {\ldots}

          SpinXY\_tauK  \textbackslash{}
/home/stafusa/ALF/pyALF/Notebooks/ALF\_data/Hubb{\ldots}  [[0.5081765927159027,
0.37862810791968693, 0.5{\ldots}

      SpinXY\_tauK\_err  \textbackslash{}
/home/stafusa/ALF/pyALF/Notebooks/ALF\_data/Hubb{\ldots}  [[0.00013766551930538375,
0.000107300494800938{\ldots}

          SpinXY\_tauR  \textbackslash{}
/home/stafusa/ALF/pyALF/Notebooks/ALF\_data/Hubb{\ldots}  [[0.0007216230187712235,
-0.07124366825946105,{\ldots}

      SpinXY\_tauR\_err  \textbackslash{}
/home/stafusa/ALF/pyALF/Notebooks/ALF\_data/Hubb{\ldots}  [[1.7167017023222387e-05,
6.709661466359632e-0{\ldots}

   SpinXY\_tau\_lattice  \textbackslash{}
/home/stafusa/ALF/pyALF/Notebooks/ALF\_data/Hubb{\ldots}  \{'L1': [4.0, 0.0], 'L2':
[0.0, 4.0], 'a1': [1{\ldots}

           SpinZ\_tauK  \textbackslash{}
/home/stafusa/ALF/pyALF/Notebooks/ALF\_data/Hubb{\ldots}  [[0.508040649079463,
0.37881927179713126, 0.50{\ldots}

       SpinZ\_tauK\_err  \textbackslash{}
/home/stafusa/ALF/pyALF/Notebooks/ALF\_data/Hubb{\ldots}  [[0.0017636045356165923,
0.0002240298516522355{\ldots}

           SpinZ\_tauR  \textbackslash{}
/home/stafusa/ALF/pyALF/Notebooks/ALF\_data/Hubb{\ldots}  [[0.0008087478953266395,
-0.07129703130271973,{\ldots}

       SpinZ\_tauR\_err  \textbackslash{}
/home/stafusa/ALF/pyALF/Notebooks/ALF\_data/Hubb{\ldots}  [[0.0003924858556099979,
0.0001948803830615219{\ldots}

    SpinZ\_tau\_lattice
/home/stafusa/ALF/pyALF/Notebooks/ALF\_data/Hubb{\ldots}  \{'L1': [4.0, 0.0], 'L2':
[0.0, 4.0], 'a1': [1{\ldots}

[1 rows x 110 columns]
\end{Verbatim}
\end{tcolorbox}
        
    \(\bullet\) The internal energy of the system (and its error) are
accessed by:

    \begin{tcolorbox}[breakable, size=fbox, boxrule=1pt, pad at break*=1mm,colback=cellbackground, colframe=cellborder]
\prompt{In}{incolor}{9}{\boxspacing}
\begin{Verbatim}[commandchars=\\\{\}]
\PY{n}{obs}\PY{o}{.}\PY{n}{iloc}\PY{p}{[}\PY{l+m+mi}{0}\PY{p}{]}\PY{p}{[}\PY{p}{[}\PY{l+s+s1}{\PYZsq{}}\PY{l+s+s1}{Ener\PYZus{}scal0}\PY{l+s+s1}{\PYZsq{}}\PY{p}{,} \PY{l+s+s1}{\PYZsq{}}\PY{l+s+s1}{Ener\PYZus{}scal0\PYZus{}err}\PY{l+s+s1}{\PYZsq{}}\PY{p}{,} \PY{l+s+s1}{\PYZsq{}}\PY{l+s+s1}{Ener\PYZus{}scal\PYZus{}sign}\PY{l+s+s1}{\PYZsq{}}\PY{p}{,} \PY{l+s+s1}{\PYZsq{}}\PY{l+s+s1}{Ener\PYZus{}scal\PYZus{}sign\PYZus{}err}\PY{l+s+s1}{\PYZsq{}}\PY{p}{]}\PY{p}{]}
\end{Verbatim}
\end{tcolorbox}

            \begin{tcolorbox}[breakable, size=fbox, boxrule=.5pt, pad at break*=1mm, opacityfill=0]
\prompt{Out}{outcolor}{9}{\boxspacing}
\begin{Verbatim}[commandchars=\\\{\}]
Ener\_scal0           -23.600856
Ener\_scal0\_err         0.000865
Ener\_scal\_sign              1.0
Ener\_scal\_sign\_err          0.0
Name: /home/stafusa/ALF/pyALF/Notebooks/ALF\_data/Hubbard\_Square, dtype: object
\end{Verbatim}
\end{tcolorbox}
        
    \(\bullet\) The simulation can be resumed by calling \texttt{sim.run()}
again, increasing the precision of results:

    \begin{tcolorbox}[breakable, size=fbox, boxrule=1pt, pad at break*=1mm,colback=cellbackground, colframe=cellborder]
\prompt{In}{incolor}{10}{\boxspacing}
\begin{Verbatim}[commandchars=\\\{\}]
\PY{n}{sim}\PY{o}{.}\PY{n}{run}\PY{p}{(}\PY{p}{)}
\PY{n}{sim}\PY{o}{.}\PY{n}{analysis}\PY{p}{(}\PY{p}{)}
\PY{n}{obs2} \PY{o}{=} \PY{n}{sim}\PY{o}{.}\PY{n}{get\PYZus{}obs}\PY{p}{(}\PY{p}{)}
\PY{n}{obs2}\PY{o}{.}\PY{n}{iloc}\PY{p}{[}\PY{l+m+mi}{0}\PY{p}{]}\PY{p}{[}\PY{p}{[}\PY{l+s+s1}{\PYZsq{}}\PY{l+s+s1}{Ener\PYZus{}scal0}\PY{l+s+s1}{\PYZsq{}}\PY{p}{,} \PY{l+s+s1}{\PYZsq{}}\PY{l+s+s1}{Ener\PYZus{}scal0\PYZus{}err}\PY{l+s+s1}{\PYZsq{}}\PY{p}{,} \PY{l+s+s1}{\PYZsq{}}\PY{l+s+s1}{Ener\PYZus{}scal\PYZus{}sign}\PY{l+s+s1}{\PYZsq{}}\PY{p}{,} \PY{l+s+s1}{\PYZsq{}}\PY{l+s+s1}{Ener\PYZus{}scal\PYZus{}sign\PYZus{}err}\PY{l+s+s1}{\PYZsq{}}\PY{p}{]}\PY{p}{]}
\end{Verbatim}
\end{tcolorbox}

    \begin{Verbatim}[commandchars=\\\{\}]
Prepare directory "/home/stafusa/ALF/pyALF/Notebooks/ALF\_data/Hubbard\_Square"
for Monte Carlo run.
Resuming previous run.
Run /home/stafusa/ALF/ALF/Prog/ALF.out
\#\#\# Analyzing /home/stafusa/ALF/pyALF/Notebooks/ALF\_data/Hubbard\_Square \#\#\#
/home/stafusa/ALF/pyALF/Notebooks
Custom observables:
Scalar observables:
Ener\_scal
Kin\_scal
Part\_scal
Pot\_scal
Histogram observables:
Equal time observables:
Den\_eq
Green\_eq
SpinT\_eq
SpinXY\_eq
SpinZ\_eq
Time displaced observables:
Den\_tau
Green\_tau
SpinT\_tau
SpinXY\_tau
SpinZ\_tau
/home/stafusa/ALF/pyALF/Notebooks/ALF\_data/Hubbard\_Square
    \end{Verbatim}

            \begin{tcolorbox}[breakable, size=fbox, boxrule=.5pt, pad at break*=1mm, opacityfill=0]
\prompt{Out}{outcolor}{10}{\boxspacing}
\begin{Verbatim}[commandchars=\\\{\}]
Ener\_scal0           -23.60119
Ener\_scal0\_err        0.000646
Ener\_scal\_sign             1.0
Ener\_scal\_sign\_err         0.0
Name: /home/stafusa/ALF/pyALF/Notebooks/ALF\_data/Hubbard\_Square, dtype: object
\end{Verbatim}
\end{tcolorbox}
        
    \begin{tcolorbox}[breakable, size=fbox, boxrule=1pt, pad at break*=1mm,colback=cellbackground, colframe=cellborder]
\prompt{In}{incolor}{11}{\boxspacing}
\begin{Verbatim}[commandchars=\\\{\}]
\PY{n+nb}{print}\PY{p}{(}\PY{l+s+s2}{\PYZdq{}}\PY{l+s+se}{\PYZbs{}n}\PY{l+s+s2}{Running again reduced the error from}\PY{l+s+se}{\PYZbs{}n}\PY{l+s+s2}{\PYZdq{}}\PY{p}{,} \PY{n}{obs}\PY{o}{.}\PY{n}{iloc}\PY{p}{[}\PY{l+m+mi}{0}\PY{p}{]}\PY{p}{[}\PY{p}{[}\PY{l+s+s1}{\PYZsq{}}\PY{l+s+s1}{Ener\PYZus{}scal0\PYZus{}err}\PY{l+s+s1}{\PYZsq{}}\PY{p}{]}\PY{p}{]}\PY{p}{,} \PY{l+s+s2}{\PYZdq{}}\PY{l+s+se}{\PYZbs{}n}\PY{l+s+s2}{to}\PY{l+s+se}{\PYZbs{}n}\PY{l+s+s2}{\PYZdq{}}\PY{p}{,} \PY{n}{obs2}\PY{o}{.}\PY{n}{iloc}\PY{p}{[}\PY{l+m+mi}{0}\PY{p}{]}\PY{p}{[}\PY{p}{[}\PY{l+s+s1}{\PYZsq{}}\PY{l+s+s1}{Ener\PYZus{}scal0\PYZus{}err}\PY{l+s+s1}{\PYZsq{}}\PY{p}{]}\PY{p}{]}\PY{p}{)}
\end{Verbatim}
\end{tcolorbox}

    \begin{Verbatim}[commandchars=\\\{\}]

Running again reduced the error from
 Ener\_scal0\_err    0.000865
Name: /home/stafusa/ALF/pyALF/Notebooks/ALF\_data/Hubbard\_Square, dtype: object
to
 Ener\_scal0\_err    0.000646
Name: /home/stafusa/ALF/pyALF/Notebooks/ALF\_data/Hubbard\_Square, dtype: object
    \end{Verbatim}

    \textbf{Note}: To run a fresh simulation - instead of performing a
refinement over previous run(s) - the Monte Carlo run directory should
be deleted before rerunning.

    \begin{center}\rule{0.5\linewidth}{0.5pt}\end{center}

\hypertarget{exercises}{%
\subsection{Exercises}\label{exercises}}

\begin{enumerate}
\def\labelenumi{\arabic{enumi}.}
\tightlist
\item
  Check info/help commands such as \texttt{help(alf\_src)},
  \texttt{alf\_src.get\_ham\_names()}, and
  \texttt{alf\_src.get\_default\_params(\textquotesingle{}Hubbard\textquotesingle{})}.
\item
  Rerun once again and check the new improvement in precision.
\item
  Look at a few other observables (\texttt{sim.analysis()} outputs the
  names of those available).
\item
  Change the lattice size by adding, e.g.,
  \texttt{\textbackslash{}"L1\textbackslash{}":\ 4,} and
  \texttt{\textbackslash{}"L2\textbackslash{}":\ 1,} to the simulation
  parameters definitions of \texttt{sim} (step 2).
\end{enumerate}


    % Add a bibliography block to the postdoc
    
    
    
